\chapter{The Weierstrass \texorpdfstring{\(\wp\)}{wp}-Function}
\label{ch:weierstrass-wp}

The analytic theory of elliptic curves centers on the Weierstrass function
\[
  \wp(z)=\frac{1}{z^2}+
  \sum_{\lambda\in\Lambda\setminus\{0\}}
  \left(\frac{1}{(z-\lambda)^2}-\frac{1}{\lambda^2}\right).
\]
The subtraction term makes the series converge normally on compact sets away from
\(\Lambda\).  The function is even, \(\Lambda\)-periodic, and has one double pole
per fundamental parallelogram.

\begin{definition}
The Eisenstein series of weight \(k\) is
\[
  G_k(\Lambda)=\sum_{\lambda\in\Lambda\setminus\{0\}}\lambda^{-k}
\]
for \(k>2\).  Set \(g_2=60G_4\) and \(g_3=140G_6\).
\end{definition}

\begin{theorem}[The differential equation]
\label{thm:wp-ode-new}
The Weierstrass function satisfies
\[
  \wp'(z)^2=4\wp(z)^3-g_2\wp(z)-g_3.
\]
\end{theorem}

\begin{proof}[Proof sketch]
Compare Laurent expansions at \(0\).  Both sides have principal part
\(4z^{-6}\), and the constants \(g_2,g_3\) are chosen so that the remaining
negative and constant terms cancel.  The difference is then a holomorphic
elliptic function, hence constant by \cref{prop:no-holomorphic-elliptic}; its
Laurent expansion shows that constant is \(0\).
\end{proof}

\begin{corollary}
The functions \(x=\wp(z)\) and \(y=\wp'(z)\) land on the affine cubic
\[
  y^2=4x^3-g_2x-g_3.
\]
\end{corollary}

This is the exact moment where the analytic face becomes algebraic.  The final
exam's minimal-polynomial problem for \(\wp'\) over \(\C(\wp)\) is the algebraic
shadow of this differential equation.

\section*{Exercises}
\addcontentsline{toc}{section}{Exercises}

\exerciseintent{Prove evenness of \(\wp\) by reindexing the defining lattice sum.}
\begin{exercise}
Use the defining lattice sum for \(\wp\).  Substitute \(-z\) for \(z\), then
reindex the lattice terms by \(\lambda\mapsto-\lambda\) to show by calculation
that \(\wp(-z)=\wp(z)\).
\end{exercise}

\exerciseintent{Use the chain rule to derive the oddness of \(\wp'\).}
\begin{exercise}
Differentiate the identity \(\wp(-z)=\wp(z)\) with respect to \(z\).  Keep track
of the chain-rule minus sign and compute the resulting parity of \(\wp'\).
\end{exercise}

\exerciseintent{Read the Weierstrass differential equation as algebraicity over \(\C(\wp)\).}
\begin{exercise}
Assuming \cref{thm:wp-ode-new}, write an explicit quadratic polynomial
\[
  Y^2-(4\wp^3-g_2\wp-g_3)\in \C(\wp)[Y]
\]
that is satisfied by \(Y=\wp'\).  Use this calculation to show that \(\wp'\) is
algebraic of degree at most \(2\) over \(\C(\wp)\).
\end{exercise}

\exerciseintent{Extract the pole orders that will feed the genus-one Riemann--Roch calculation.}
\begin{exercise}
Using the Laurent expansions
\[
  \wp(z)=z^{-2}+O(z^2),\qquad \wp'(z)=-2z^{-3}+O(z),
\]
compute the pole orders of \(\wp\) and \(\wp'\) at \(O\in T\).
\end{exercise}

\exerciseintent{Substitute the Weierstrass functions into the cubic equation.}
\begin{exercise}
Set \(x=\wp(z)\) and \(y=\wp'(z)\).  Substitute these expressions into
\cref{thm:wp-ode-new} and verify directly that the pair \((x,y)\) lies on the
affine cubic
\[
  y^2=4x^3-g_2x-g_3.
\]
\end{exercise}
